% SPDX-License-Identifier: Apache-2.0
%
% The TxHDL runtime: the library that is the language, included verbatim.
% Built by //docs:runtime. One column, because it is mostly listings.
\documentclass[journal,onecolumn,9pt]{IEEEtran}

% One column: 80 columns fit at the body size, so nothing is reduced.
\newcommand{\listingsize}{\normalsize}
% SPDX-License-Identifier: Apache-2.0
%
% The house preamble, shared by every document in this directory. Loaded
% after \documentclass, before \title. Keeping it in one file is what
% stops two articles drifting apart in font, listing style or figure
% style.
% Computer Modern. IEEEtran selects Times (ptm) by default, so the face has
% to be chosen explicitly.
%
% Latin Modern is the Computer Modern variety used here, and the choice is
% forced rather than aesthetic. Under T1 encoding, plain `cmr` has no Type 1
% outlines unless cm-super is installed, and it is not in the pinned
% distribution, so pdflatex falls back to EC bitmaps and every glyph in the
% PDF becomes a Type 3 font. That was measured: the first build produced a
% document with 10 Type 3 fonts and no Type 1. Latin Modern is Computer
% Modern redrawn with a full T1 repertoire and real .pfb outlines, so the
% same design comes out scalable.
\usepackage[T1]{fontenc}
\usepackage[utf8]{inputenc}
\usepackage{lmodern}
% microtype is not in the pinned TeX distribution. emergencystretch alone
% keeps the two-column measure from producing overfull lines.
\setlength{\emergencystretch}{3em}

\usepackage{amsmath}
\usepackage{amssymb}
\usepackage{amsthm}
\usepackage{booktabs}
\usepackage{array}
% enumitem is not in the pinned TeX distribution; the list spacing below
% does what \setlist would have done.
\usepackage{float}
\usepackage{xcolor}
\usepackage{listings}
\usepackage{tikz}
\usetikzlibrary{arrows.meta,positioning,fit,backgrounds,calc}
\usepackage[hidelinks,bookmarks=true]{hyperref}

% Numbered, definition-style box for a rule the reader is meant to apply.
\theoremstyle{definition}
\newtheorem{rulebox}{Rule}
\newtheorem{example}{Example}

% Unnumbered asides.
\theoremstyle{remark}
\newtheorem*{tip}{Tip}
\newtheorem*{pitfall}{Pitfall}

% Tighter lists, in place of enumitem's \setlist.
\setlength{\topsep}{2pt}
\setlength{\itemsep}{1pt}
\setlength{\parsep}{0pt}

\newcommand{\code}[1]{\texttt{#1}}
\newcommand{\term}[1]{\emph{#1}}

% Syntax highlighting. Four roles, distinguishable in colour and still
% distinguishable in greyscale, because an IEEE article gets printed: the
% keyword is the only bold role, the comment is the only italic one, and the
% two remaining roles differ in lightness as well as in hue.
\definecolor{kwblue}{RGB}{20,60,150}
\definecolor{cmtgreen}{RGB}{90,120,90}
\definecolor{strred}{RGB}{150,50,40}
\definecolor{numgray}{gray}{0.45}
\definecolor{framegray}{gray}{0.70}
\definecolor{bgshade}{gray}{0.97}
% A darker fill for figure cells. The listing background has to stay very
% light to keep code readable; a figure cell has to be clearly shaded or the
% distinction it draws is invisible in print.
\definecolor{cellshade}{gray}{0.90}

% The listing size is the document's choice, because it depends on the
% column: a two-column article sets `\listingsize` to 6pt and keeps its
% listings within 70 columns, a one-column document keeps the body size
% and includes source files at 80. Lines are never broken; a line that does not fit
% overflows the frame, which is visible, rather than wrapping, which is
% not.
\providecommand{\listingsize}{\scriptsize}

% `'` is deliberately not a string delimiter: the language uses it for loop
% labels, as Rust does, so treating it as a quote would swallow the rest of
% the line.
\lstdefinestyle{txhdl}{
  basicstyle=\ttfamily\listingsize,
  keywordstyle=\bfseries\color{kwblue},
  commentstyle=\itshape\color{cmtgreen},
  stringstyle=\color{strred},
  numberstyle=\tiny\color{numgray},
  morekeywords={mod,use,pub,const,type,struct,enum,trait,impl,fn,async,let,mut,
    match,if,else,loop,while,for,in,break,continue,return,as,await,
    module,bus,channel,transaction,protocol,pipeline,flow,seq,config,
    port,stage,pipe,instance,connect,bind,tag,domain,blackbox,foreign,
    property,role,signal,handler,true,false,
    sequence,interface,modport,implements,package,import,var,wire,func,proc,
    case,default,next,fork,branch,join,state,fsm},
  morecomment=[l]{//},
  morestring=[b]",
  showstringspaces=false,
  columns=fullflexible,
  keepspaces=true,
  breaklines=false,
  % Framed, so a listing is bounded rather than floating in the column.
  frame=single,
  rulecolor=\color{framegray},
  framesep=3pt,
  backgroundcolor=\color{bgshade},
  xleftmargin=3pt,
  xrightmargin=3pt,
  aboveskip=6pt,
  belowskip=6pt,
  % Every listing is numbered and captioned, so it can be referred to by
  % number. `caption` without `float` numbers the listing and leaves it
  % where it was written, which is what a listing in running prose needs.
  captionpos=b,
  abovecaptionskip=2pt,
  belowcaptionskip=6pt,
}

% Figures drawn as text. Framed the same way, and with the highlighting
% switched off, because a box-and-arrow diagram is not code and half its
% words turning blue reads as an accident. Setting a style adds keys rather
% than clearing the ones already set, so the three roles are neutralised by
% hand rather than by leaving them out.
\lstdefinestyle{ascii}{
  basicstyle=\ttfamily\listingsize,
  keywordstyle=\ttfamily,
  commentstyle=\ttfamily,
  stringstyle=\ttfamily,
  columns=fullflexible,
  keepspaces=true,
  frame=single,
  rulecolor=\color{framegray},
  framesep=3pt,
  backgroundcolor=\color{bgshade},
  xleftmargin=3pt,
  xrightmargin=3pt,
  aboveskip=6pt,
  belowskip=6pt,
}
\lstset{style=txhdl}

% House TikZ styles. `shorten` lives on the arrow styles rather than on each
% arrow, so every arrow in the document gets the gap, including ones added
% later. TikZ clips a path at a node's border, so without it an arrowhead
% butts against the box with no gap at all. Keep `node distance` well above
% twice the shorten, or the arrow shrinks to nothing.
\tikzset{
  box/.style={draw,rounded corners=2pt,align=center,inner sep=4pt,
              font=\footnotesize},
  boxf/.style={box,fill=bgshade},
  lbl/.style={font=\scriptsize\itshape,align=center},
  fl/.style={-{Stealth[length=4pt]},shorten >=2.5pt,shorten <=2.5pt},
  fld/.style={fl,dashed},
}



\InputIfFileExists{buildstamp.tex}{}{}
\providecommand{\buildstamp}{Draft build (outside the Bazel flow).}

\title{The TxHDL Runtime}

\author{Filip Filmar and Dragiša Janković%
\thanks{Both authors are independent. Filip Filmar is the
corresponding author, at f@hdlfactory.com; Dragiša Janković is
reached at linkedin.com/in/dragisa-jankovic.}%
\thanks{This document holds the runtime library of TxHDL, the crate
\code{txhdl} at \code{//lib} of the project repository, included
verbatim from the tree at build time. The language it implements is
described in the embedding article~\cite{embedding}; the examples
written against it are a third document~\cite{examples}; the cover
document~\cite{cover} lists all of them.}%
\thanks{The authors used a large language model, Claude, as an
assistant in exploring the concepts and in writing and constructing
the documents and the programs; every commit of the repository says
so and carries the prompts that led to it, verbatim.}%
\thanks{\buildstamp}}

\markboth{The TxHDL Runtime}{Filmar and Janković: The TxHDL Runtime}

\begin{document}
\maketitle

\begin{abstract}
TxHDL is a hardware description language embedded in Rust, and once
every construct Rust provides has been deleted from it, what is left is
a library. This document is that library. It states the model of time
the executor keeps, which is the one part of the runtime a reader
cannot recover from the signatures alone, and then includes each source
file of the crate in full. Nothing here has been lowered to hardware;
the prototype simulates, and it is what the lowering experiment would
be written against.
\end{abstract}

\begin{IEEEkeywords}
Hardware description languages, embedded domain-specific languages, Rust,
runtime, simulation.
\end{IEEEkeywords}

{\footnotesize\tableofcontents}

% SPDX-License-Identifier: Apache-2.0
\section{What the Runtime Is}
\label{sec:r-overview}

The embedding article~\cite{embedding} argues that a hardware
description language embedded in Rust needs no parser and no keywords,
and that what remains is a library. The crate \code{txhdl} is that
library, and this document is the crate. Every file below is included
from the tree when the document is built, so a listing here cannot
differ from the code, and the rule the project keeps is that a concept
enters the runtime, a compiled example and these documents in the same
change, or the change is not finished.

The crate is small on purpose. Four modules and one procedural macro
crate hold everything:

\begin{itemize}
\item \code{txhdl::types}: two-valued \code{Bit}, nine-valued
  \code{Logic}, the numeric vectors \code{U<N>} and \code{I<N>},
  \code{logic::Vec<N>}, and the two marker traits \code{Transaction}
  and \code{Tag}.
\item \code{txhdl::comp}: clocks, wires and channels and their ends,
  interfaces, crossings, registers and memories, units, configurations,
  and the executor.
\item \code{txhdl::pipeline}: operators that take a cycle, all
  \code{async}.
\item \code{txhdl::funcs}: operators that cannot, all plain \code{fn}.
\item \code{txhdl\_macros}: the two derives, \code{interface!},
  \code{when!} and \code{case!}, without \code{syn} or \code{quote}.
\end{itemize}

The whole crate is a simulation-shaped prototype: values exist while it
runs, wires are shared cells, and a testbench may look at anything.
The lowering to hardware is the experiment the article states, and this
crate is what that experiment would be written against.

\subsection{The model of time}
\label{sec:r-time}

The executor is the one part of the runtime a reader cannot recover
from the signatures, so it is stated here before the code.

\textbf{Time runs in ticks.} One poll of the top unit is one tick.
Every clock is a type with a period and a phase in ticks, and has a
rising edge at every tick \code{t} with \code{t >= PHASE} and
\code{(t - PHASE) \% PERIOD == 0}, and a falling edge half a period
later. The default clock has period 2 and phase 0, so its cycle is two
ticks and its falling edge is a tick of its own; a single-clock design
mentions neither number, and its testbench steps by
\code{Running::cycle}.

\textbf{A process waits, then reads.} The waits are events:
\code{C::rising()} and \code{C::falling()}, the next edge of a clock;
\code{rx.wait()}, the next rising edge of the channel's clock at which
a transaction is offered, taken; and \code{until(C::rising, || cond)},
or the same with \code{falling}, the next such edge at which a
condition on state holds. Reads are plain: \code{Reg::get} returns
what the last edge latched, a wire's \code{get} returns what it
carries, a memory's \code{read} is its asynchronous port. Drives are
deferred: \code{Reg::set} and \code{Mem::write} take effect at the end
of the step, so a read after a drive in the same iteration still sees
what the edge latched.

\textbf{Every \code{.await} is a cycle boundary.} A process that has
crossed an edge at this step cannot cross it again, so a second wait
written after a first waits for the next edge, as an operator does.
Waits that share one edge are written inside \code{parallel!}, which
polls its futures together and returns their outputs as a tuple; that
is the one case the executor treats specially. Once one wait of the
group has crossed the edge at this step, the others in the group are
free, each once, so the group is one wait. State read before a wait,
at the top of a loop, is read in the previous step and before that
step's drives commit; state is read after the wait.

\textbf{A process is its waker.} The executor tells processes apart by
the data pointer of the waker they are polled with. \code{join2} and
\code{join\_all} give each child a fresh one, so the rule above holds
per process rather than per unit. \code{parallel!} keeps its futures in
the one process, which is the difference between the two: \code{join2}
forks, \code{parallel!} does not.

\textbf{An operator's cycle is in the process's clock.} An operator in
\code{txhdl::pipeline} waits one edge of whatever clock the process
last crossed an edge of, which is the only sense a latency has inside
a process. A process that has crossed no edge yet is in the default
clock.

\textbf{What the executor does not check.} A process that reads a
register of another clock without a \code{Crossing} is a hazard, and
the prototype runs it. The register's type says its clock; the
process has no type that says one. The asynchronous FIFO example
commits this hazard on purpose, on its memory, under the pointer
discipline that makes it safe.

\subsection{What is prototype only}

Two things in the code exist for the simulation and would not be
lowered. \code{Running} and \code{simulate} drive a design from a
\code{main}, and \code{now} reports the time step for traces. And the
shared cells behind every wire, channel and register, with the deferred
drives that commit at the end of a step, are how a simulation shares
state; the lowering would replace them with nets and flops.

\subsection{Tracing}

A design can be watched as a waveform. \code{comp::trace} holds a
registry of probes, each a dotted path, a width and a closure that
reads the signal; \code{\#[derive(Trace)]} on a unit registers every
field under its field name, so the hierarchy is the nesting of units
and nothing is named at construction. Registers, wire ends, channel
ends and crossings register themselves, plain values register nothing,
and \code{\#[derive(Value)]} gives a struct or a fieldless enum of the
design's own its width and bits. \code{Vcd} is the sink: the executor
calls it at the end of every tick, after the drives have committed,
and it writes what changed; a clock is written as the level it is,
high from its rising edge to its falling one. FST would be the same probes with another
writer, and there is one now: \code{Fst}, over the \code{fst-writer}
crate, the project's first external dependency, pinned by lockfile
through \code{crate\_universe}. \code{Wave::from\_env} chooses the sink
by the environment, \code{TXHDL\_FST} or \code{TXHDL\_VCD}, and
\code{stop} finishes the file. A compound value is one bit vector and
one signal per field beside it, from \code{Value::parts}; the writer
has no string variables, so a value keeps its bits.
The examples document draws a VCD through \code{vcdcvt} and
\code{sqlite2drawtiming}, prebuilt tools pinned by checksum, and
\code{//tools/dt2tikz}, which turns drawtiming's text into TikZ.

\subsection{Netlists}

The probe registry records, for every signal, its kind and the shared
cell it is on, so the two ends of one wire match. \code{netlist::verilog}
walks a design through it and writes a Verilog skeleton: a module per
unit, a \code{reg} per register, ports where a unit holds one end of
a wire whose other end is elsewhere, and instances with their ports
connected. It sees fields only, so a unit whose ports are parameters
of \code{run} yields its registers and nothing else, and it writes no
behaviour. That is the structural half of the lowering. The
behavioural half has begun: \code{\#[pipeline]} lowers an \code{async
fn} of awaited operators, and \code{\#[lower]} lowers a unit's
\code{run} loop, with \code{netlist::unit\_verilog} assembling the
module from the unit's fields, the ports of \code{run} and the
statements the macro read, \code{case!} and channels included. The
macro builds an expression tree, \code{netlist::Expr}, and a
\code{Lowered} unit around it, and two emitters render it, Verilog
and VHDL-2008. The VHDL is what closes the loop: the build simulates
a lowered unit under \code{nvc}, with a testbench
\code{//tools/fst2tb} generates from the trace the Rust run wrote, and
the test passes only if registers and outputs agree at every tick.
What neither emitter reads yet is a unit's second process, a
falling-edge wait, and a memory; and a unit with channels cannot be
checked until the runtime's channel is a cycle-accurate handshake
rather than a mailbox.

\subsection{Building it}

\begin{lstlisting}[style=ascii,caption={The runtime and everything written against it, built hermetically.},label={lst:r-build}]
bazel build //lib/...                 # the crate and the macros
bazel build //lib/examples/...        # every example
bazel run //lib/examples:ex_fifo      # one that prints a trace
bazel run @rules_rust//:rustfmt \
  --@rules_rust//:rustfmt.toml=//:rustfmt.toml -- //lib/...
\end{lstlisting}

The last line formats the tree at 80 columns, which is the width these
documents include a file at. A line longer than that overflows its
frame, visibly, rather than wrapping.

% SPDX-License-Identifier: Apache-2.0
\section{\code{txhdl::types}}
\label{sec:r-types}

Values, and the two markers. \code{Bit} is two valued, \code{Logic} is
IEEE 1164's nine, and \code{U<N>} and \code{I<N>} are the numeric
vectors, with the width a const parameter. A width that is an
expression of other widths, \code{U<\{A + B\}>}, needs nightly Rust,
which is the largest constraint the embedding found, and the
prototype declares each width by hand instead. The conversions from
Rust integers are what let a literal be written where a value is
expected, and the narrowing ones are checked in debug builds.

\code{Transaction} marks a struct that moves over a channel, and a
bare word is one too. \code{Tag} is a synchronisation domain and its
policy, as a type; which clock edge moves the data is a
\code{Clock}, in the next module.

\lstinputlisting[caption={\code{lib/src/types.rs}.},label={lst:r-types}]{lib/src/types.rs}

\section{\code{txhdl::netlist}}
\label{sec:r-netlist}

The lowering's target, Section~\ref{sec:r-lowering}: \code{Expr},
\code{Target} and \code{Stmt}, which \code{\#[lower]} builds; the
\code{Lowered} unit that holds them with its fields, ports, wires and
memory contents; the two emitters, \code{verilog} and \code{vhdl},
with the width inference an extension needs and the hoisting of a
slice of an expression into a wire; the ports sidecar a testbench
generator reads; and, first in the file, the structural skeleton from
the trace walk, which sees fields only.

\lstinputlisting[caption={\code{lib/src/netlist.rs}.},label={lst:r-netlist}]{lib/src/netlist.rs}

% SPDX-License-Identifier: Apache-2.0
\section{\code{txhdl::comp}}
\label{sec:r-comp}

The composition module: everything a unit is made of, and the executor
that runs it. In order of appearance: clocks with their period and
phase; a wire and a channel, each split into a driving end that is not
\code{Clone} and a reading end that is; the \code{Member} trait that
lets \code{interface!} split a member without knowing which it is;
\code{Crossing}, the only way between two clocks; registers and
memories; \code{mux}; the \code{Module} trait, whose \code{run} takes
the unit's inputs and outputs as type parameters; configurations and
the \code{config!} macro; \code{join2} and \code{join\_all}, which give
each child its own process identity; and the executor, whose rules
Section~\ref{sec:r-time} states.

\lstinputlisting[caption={\code{lib/src/comp.rs}.},label={lst:r-comp}]{lib/src/comp.rs}

% SPDX-License-Identifier: Apache-2.0
\section{\code{txhdl::pipeline} and \code{txhdl::funcs}}
\label{sec:r-operators}

Two modules that hold the same kind of thing, operators, and differ in
one property: whether the operator takes a cycle. An operator in
\code{pipeline} is \code{async} and awaits one edge of the process's
clock, so a pipeline's depth follows from the operators it awaits and
nobody places a stage boundary. An operator in \code{funcs} is a plain
\code{fn}, has no \code{.await} to write, and cannot count cycles, which
is the whole enforcement of the rule that a combinational body is
timeless. The operators with a Rust spelling are Rust's, on the value
types of Section~\ref{sec:r-types}; \code{funcs} holds the two that
have none, the arithmetic shift and the signed compare, and the
halves of a word.

\code{drive} is what makes an \code{async fn} a pipeline rather than
a sequence: at every edge at which an input is offered it starts an
invocation, it polls every invocation in flight as a process of its
own, each started as having crossed the edge it began at so its first
wait is the next edge, and it sends each result as its invocation
completes, oldest first. Several invocations are in flight at once,
each at a different await.

\lstinputlisting[caption={\code{lib/src/pipeline.rs}.},label={lst:r-pipeline}]{lib/src/pipeline.rs}
\lstinputlisting[caption={\code{lib/src/funcs.rs}.},label={lst:r-funcs}]{lib/src/funcs.rs}

% SPDX-License-Identifier: Apache-2.0
\section{\code{txhdl\_macros}}
\label{sec:r-macros}

The procedural macros, written against \code{proc\_macro} alone so the
crate has no dependencies. \code{\#[derive(Transaction)]} and
\code{\#[derive(Bus)]} emit one empty \code{impl} each.
\code{\#[derive(Value)]} gives a struct of values, or a fieldless
enum, its width and its bits for a trace, and \code{\#[derive(Trace)]}
registers every field of a unit under its field name; both read the
item's generic parameters so a generic unit can derive them.
\code{\#[pipeline(op = latency, ..)]} is the first of the lowering: on
an \code{async fn} of awaited operator calls it keeps the function and
emits a \code{NAME\_verilog()} beside it, each operator a stage of the
stated latency and each value delayed to meet the operator that uses
it. Its grammar is the one \code{mac} needs and no wider, and it
refuses the rest with its own message. \code{\#[lower]} on an
\code{impl Unit for Unit} does the same for a unit with state,
filling the header's \code{<In, Out>} in from \code{run}'s signature,
so the ports are named once: it reads \code{run} as one loop, or
several under
\code{join2}, each a process that waits for the rising or the falling
edge of its clock, reads registers and ports into names, predicates
drives with
\code{with!}, \code{case!} and \code{if}, chooses values
with
\code{select!},
and drives registers, words of memories and output ports, and writes
\code{Unit::lowered(name)}, \code{verilog} and \code{vhdl} beside
the impl. Ports come from \code{run}'s signature, a side of several
written as a tuple; registers, memories and widths come from the
struct through the \code{Trace} derive, which also derives
\code{Fields} and \code{Port}; and the configuration's constants are
evaluated when \code{lowered} runs, so a generic unit lowers once per
build. A \code{let} that computes is a wire named for it, a read of a
port or register an alias. It reads \code{case!} and \code{select!}
as chains on the value, with an enum variant, an integer or \code{\_}
as a pattern and a guard joined by \code{\&\&}; \code{slice},
\code{concat}, \code{sext} and \code{zext} by their turbofish, every
width it needs written, since it sees tokens and not types, and only
\code{concat}'s low operand left to Rust as \code{\_}; a
struct literal as the concatenation of its fields; an \code{if} chain
as arms under their conditions, each arm's statements lowered in
turn, a \code{send} in an arm hoisted with the arm's path as
\code{valid}, and a \code{set} of an output under it refused, since
an output is a wire; \code{self.m.at(a)}
on the left of a drive as a word of a memory and \code{self.m.read(a)}
as an index; and channels as a valid/ready handshake: \code{until} or
a channel's \code{wait} is the guard on every register drive after
it, a receive asserts \code{ready} under the guard, a \code{take}
binds \code{valid} and \code{data} to its two names and asserts
\code{ready} the same way, a send drives
\code{data} and asserts \code{valid} under it. What it does not read
it refuses with a message that names the construct.
A field of a port's value, \code{p.tag} where \code{p} is what a
channel port carries, is read as the slice of the port's data the
field occupies: the \code{Value} derive lays a struct's fields out
first field highest and says so in a \code{layout}, and the netlist
finds the bits when \code{lowered} runs, so a generic field's width
is no obstacle. \code{station!(Name, N)} writes a reservation
station of \code{N} inputs as text, a line struct, the station's
struct and its lowered \code{run}, every wire of the step named for
what it is, and \code{TXHDL\_MACRO\_DUMP} writes the text out for
reading; the attribute in its output expands as any does.
\code{interface!} takes an interface and any number of roles, checks
that every role names every member once and that no member is driven
twice, and emits a struct per role plus a constructor that splits each
member. \code{with!(self <= \{ .. \})} is the drives of one struct,
its name written once: an entry \code{field: value} is a drive,
\code{c ? field: value} a drive under a condition, \code{c ? \{ ..
\} else \{ .. \}} a group under one with the entries after
\code{else} under its failure, and a memory word \code{m.at(i)} a
field like any other; entries apply in order, the last drive of a
field winning, and each lowers as an \code{if} in the clocked block;
an entry that is not \code{field: value} is refused with the macro's
own message. \code{when!(c => self \{ .. \} else \{ .. \})} is one
group of it with the condition written first, the same entries and
the same \code{if}; an \code{if} whose condition is a constant of
the build is folded, its arm kept or dropped and the test gone. \code{case!} is a predicated drive per arm that reads
as a conditional: each statement inside an arm is \code{register <=
value}, the arrow pointing into the register, and a statement that
is not is refused with the macro's own message. Both are procedural
because a \code{macro\_rules!} \code{expr} fragment may be followed
only by \code{=>}, \code{,} or \code{;}, so \code{lhs <= rhs} cannot be
matched declaratively at all.

\lstinputlisting[caption={\code{lib/macros/lib.rs}.},label={lst:r-macros}]{lib/macros/lib.rs}

% SPDX-License-Identifier: Apache-2.0
\section{\code{txhdl\_parts}}
\label{sec:r-parts}

What the language provides as hardware, written in its own lowerable
subset and checked as every unit is. A unit's channel port is one
side of a channel: the lowering makes of it a valid, a ready and the
data, and the testbench that checks the unit plays the channel's
other side from the trace. The channel itself, the elastic buffer of
two between two units, is not in either unit, so two lowered units
wired port to port would have no buffer between them, and their
handshakes would form a combinational loop, which a placer refuses.
\code{Buffer<W>} is the channel as hardware, one per channel of a run
wherever two lowered units meet, on a board or in a larger netlist.

Its rules are the runtime's, from Section~\ref{sec:r-comp}, with no
freedom taken. A take moves the tail into the head. A push goes into
the head if it is empty after that, else into the tail. \code{ready}
on the sending side is the tail being empty as the edge left it,
\code{valid} on the receiving side the head being full, and the data
is the head; all three are registers, so two units on a channel have
no combinational path between them. That the part is the runtime's
channel bit for bit is checked by the examples
document~\cite{examples}, which drives one with a pattern of offers
and takes, does the same on a runtime channel, and asserts after
every cycle that the two agree; the part's netlist is then checked
against that trace under both simulators.

\lstinputlisting[caption={\code{lib/parts/src/buffer.rs}.},label={lst:r-buffer}]{lib/parts/src/buffer.rs}

% SPDX-License-Identifier: Apache-2.0
\section{\code{txhdl::foreign}}
\label{sec:r-foreign}

A module written in Verilog as a unit. Verilator compiles the module
into a C++ model, and \code{//tools/vshim} writes two things from the
module's port list: a C shim over the model, and the Rust that makes
a \code{Unit} of it. The shim is four calls and a handle: settle the
combinational logic on the inputs as they stand, pulse the clock,
set an input port by index, read a port by index, values as four
words of thirty-two bits, so a port up to 128 bits wide crosses.
\code{Model} holds the handle and the shim's functions and frees the
model when dropped. The generated unit has an \code{In} or
\code{Out} per port and an \code{Rx} or \code{Tx} per \code{x\_data},
\code{x\_valid}, \code{x\_ready} trio, and its step, once per rising
edge, is: the inputs in as the edge left them, a channel's head and
its offer, and the room the runtime has for the module's sends; a
settle; the transfers the module agrees to, a take where it says
ready and an offer stands, a send where it says valid and there is
room; the edge; and the outputs out. So the module is a registered
unit, its outputs seen a step later, and goes first in the join
order as the timer of the core does. What it cannot do is lower: the
lowering carries it as nothing, and a netlist that holds one is the
board top's business, by hand, as the core's is.
\code{verilog\_unit()} in \code{lib/foreign.bzl} builds the model,
the shim and the crate from the one Verilog file.

A VHDL entity goes the same way with nvc as the engine, run as a
child process, since nvc is a program and not a library. The tool
writes a testbench around the entity that reads a line per step
from its standard input, a command and the inputs as hex, applies
them, settles or pulses the clock, and writes the outputs back as
one line; the build analyses and elaborates the two and makes the
script that runs nvc on them, an executable of the build's own,
since a test target cannot be depended on; and \code{Cosim} starts
the script in the runfiles, where its paths hold, and speaks the
lines, a settle and an edge each one exchange. The same generated
unit drives either engine, since both answer to set, get, settle
and edge; a step costs a pipe's round trip under nvc, some tens of
microseconds, and nothing under Verilator. \code{vhdl\_unit()}
builds the testbench, the script and the crate. The examples
document~\cite{examples} co-runs a unit with its own Verilog and
its own VHDL at once.

\lstinputlisting[caption={\code{lib/src/foreign.rs}.},label={lst:r-foreign}]{lib/src/foreign.rs}

% SPDX-License-Identifier: Apache-2.0
\section{\code{txhdl::regmap}}
\label{sec:r-regmap}

What a register map declared with \code{regmap!} stands for at run
time. The declaration is the macro: a peripheral says its map once, a
word index, a name, an access and a sentence a register, and under a
register its fields, each a name, a low bit, a width, an access, a
reset value and a sentence. From that the macro writes a module of
constants a program addresses by, the byte offset of each register
and a \code{Field} for each field, the read mux and the write enables
a unit's \code{run} calls, a function a field that takes it out of a
written word and one a register that packs its fields into the word,
and a table. The lowering reads the same declaration in the file and
inlines the same functions into the netlist, so the map a program sees
and the map the hardware decodes are one text; the examples
document~\cite{examples} shows a peripheral written that way.

This module is the data. \code{Field} is the constant a program uses:
its shift, its width and its reset, with \code{mask}, \code{get},
\code{set} and \code{with}, all of them \code{const}, so a field's
mask is a constant where a program wants one. \code{RegMap} is the
table, its \code{Reg}s and their \code{FieldInfo}s, which
\code{//tools/regmap} writes as a C header, a define a register with
its offset and four a field, or as a device tree node at a base, and
\code{//tools/datasheet} writes as a sheet's register table, a row a
register and a row a field under it.

\lstinputlisting[caption={\code{lib/src/regmap.rs}.},label={lst:r-regmap}]{lib/src/regmap.rs}


\begin{thebibliography}{10}

\bibitem{embedding}
F.~Filmar and D.~Janković, ``Deleting a language: Embedding TxHDL in
Rust,'' project repository \code{HDL/txhdl}, \code{//docs:embedding},
2026.

\bibitem{examples}
F.~Filmar and D.~Janković, ``TxHDL by example,'' project repository
\code{HDL/txhdl}, \code{//docs:examples}, 2026.

\bibitem{cover}
F.~Filmar and D.~Janković, ``TxHDL documents,'' project repository
\code{HDL/txhdl}, \code{//docs:cover}, 2026.

\bibitem{merge}
F.~Filmar and D.~Janković, ``TxHDL: Unifying a dataflow and a transaction
hardware description language,'' project repository \code{HDL/txhdl},
\code{//docs:article}, 2026.

\end{thebibliography}

\end{document}
